\documentclass[aspectratio=169,13pt,t,fleqn] {beamer}
%
% Choose how your presentation looks.
%
% For more themes, color themes and font themes, see:
% http://deic.uab.es/~iblanes/beamer_gallery/index_by_theme.html
%
\mode<presentation>
{
  \usetheme{Madrid}      % or try Darmstadt, Madrid, Warsaw, ...
  \usecolortheme{beaver} % or try albatross, beaver, crane, ...
  \usefonttheme{serif}  % or try serif, structurebold, ...
%  \setbeamertemplate{navigation symbols}{}
  \setbeamertemplate{caption}[numbered]
  \setbeamertemplate{items}[default]
}

\usepackage[,UTF8,linespread=1.2]{ctex}
\usepackage{wrapfig}
\usepackage{caption}
\captionsetup[figure]{name=图,labelsep=period}
\usepackage[,most]{tcolorbox}
\usepackage{picinpar}


\usepackage{xcolor}
\definecolor{xy_blue}{cmyk}{0.6,0.35,0.0,0.4}
\definecolor{xy_orange}{cmyk}{0.0,0.7,1.0,0.0}
\definecolor{xy_green}{cmyk}{0.2,0.0,1.0,0.15}
\xdefinecolor{xy_gold}{RGB}{135,103,79}
\xdefinecolor{xy2_gold}{cmyk}{0,0.35,0.75,0.05}


\setbeamercolor{section number projected}{fg=red,bg=red!20}%控制球的颜色
\setbeamercolor{subsection number projected}{fg=blue,bg=blue!40}
\setbeamercolor{item}{fg=xy_gold} %设置列表项颜色

\usepackage{amsthm}
\newtheorem{peroration}{结论：}
\newtheorem{exploration}{探究：}
\newtheorem{prooft}{证明：}
\newtheorem{question}{问题}
\newtheorem*{solve}{解:}
\newtheorem{instance}{例}   % 解：独立计数器
\BeforeBeginEnvironment{peroration}{%
	\setbeamercolor{block title}{bg=xy_green,fg=white}
	\setbeamercolor{block body}{bg=xy_green!10}
}
\AfterEndEnvironment{peroration}{
	\setbeamercolor{block title}{bg=xy_blue,fg=white}
	\setbeamercolor{block body}{bg=xy_blue!10}
}
\BeforeBeginEnvironment{exploration}{%
	\setbeamercolor{block title}{bg=xy_orange,fg=white}
	\setbeamercolor{block body}{bg=xy_orange!10}
}
\AfterEndEnvironment{peroration}{
	\setbeamercolor{block title}{bg=xy_blue,fg=white}
	\setbeamercolor{block body}{bg=xy_blue!10}
}
\BeforeBeginEnvironment{solution}{%
	\setbeamercolor{block title}{bg=xy_green,fg=white}
	\setbeamercolor{block body}{bg=xy_green!10}
}
\AfterEndEnvironment{solution}{
	\setbeamercolor{block title}{bg=xy_green,fg=white}
	\setbeamercolor{block body}{bg=xy_green!10}
}
\BeforeBeginEnvironment{prooft}{%
	\setbeamercolor{block title}{bg=xy_green,fg=white}
	\setbeamercolor{block body}{bg=xy_green!10}
}
\AfterEndEnvironment{prooft}{
	\setbeamercolor{block title}{bg=xy_green,fg=white}
	\setbeamercolor{block body}{bg=xy_green!10}
}
\BeforeBeginEnvironment{question}{%
	\setbeamercolor{block title}{bg=xy_orange,fg=white}
	\setbeamercolor{block body}{bg=xy_orange!10}
}
\AfterEndEnvironment{question}{
	\setbeamercolor{block title}{bg=xy_green,fg=white}
	\setbeamercolor{block body}{bg=xy_green!10}
}
\BeforeBeginEnvironment{solve}{%解：
	\begingroup
	\setbeamercolor{block title}{bg=xy_orange!60,fg=white}
	\setbeamercolor{block body}{bg=xy_orange!10}
}
\AfterEndEnvironment{solve}{%解：
	\endgroup
}

\setbeamercolor{prooft}{fg=red, bg=green!10}       % 标题颜色
\setbeamercolor{prooft body}{fg=red, bg=red!5}  % 内容块背景
\setbeamerfont{prooft head}{series=\bfseries}       % 标题字体
\setbeamerfont{question body}{shape=\upshape}         % 内容字体
\renewcommand{\thefigure}{\arabic{figure}}

\newtcolorbox{summarybox}[1][]{
	colback=purple!10,
	colframe=purple!70!black,
	fonttitle=\bfseries,
	title=#1,
	arc=1mm,
	boxrule=0pt,drop shadow,
	left=2mm, right=2mm, top=2mm, bottom=2mm,
}
\newtcolorbox{summarybox-answer}{
	width=0.9\textwidth,
	colback=white,
	blanker,
	borderline west={2pt}{20pt}{cyan},
	left=10mm, right=2mm, top=2mm, bottom=2mm,
	}
\title[4.1指数]{4.1\ 指数}
\subtitle{n次方根与分数指数幂、无理数指数幂}
\author[好好学习，努力向上]{}
\institute{}
\date{2026.11.5}
%\setbeamertemplate{theorems}[numbered]
\theoremstyle{definition}
\newtheorem{exam}{例}
\AtBeginSection[]
{
  \begin{frame}<beamer>
    \frametitle{本节主要内容}
    \tableofcontents[currentsection,currentsubsection]
  \end{frame}
}

\AtBeginDocument{
	\setlength{\abovedisplayskip}{3pt}
	\setlength{\belowdisplayskip}{3pt}
	\setbeamertemplate{theorems}[numbered]
}

\usepackage{pgfplots}
\pgfplotsset{compat=1.18}
\usetikzlibrary{arrows.meta}

\begin{document}
\begin{frame}{}
	\titlepage
\end{frame}

\begin{frame}{内容概览}
	\tableofcontents
\end{frame}

\section{n次方根}

\begin{frame}{回顾：平方根与立方根}
	我们已经知道：
	\begin{itemize}
		\item 如果 \(x^2=a\)，那么 \(x\) 叫做 \(a\) 的\textbf{平方根}. 例如，\(\pm 2\) 就是 \(4\) 的平方根.
		\item 如果 \(x^3=a\)，那么 \(x\) 叫做 \(a\) 的\textbf{立方根}. 例如，\(2\) 就是 \(8\) 的立方根.
	\end{itemize}
	\vspace{1em}
	\pause
	类似地，由于 \((\pm 2)^4=16\)，我们把 \(\pm 2\) 叫做 \(16\) 的 \(4\) 次方根；由于 \(2^5=32\)，\(2\) 叫做 \(32\) 的 \(5\) 次方根.
\end{frame}

\begin{frame}{n次方根的定义}
	\begin{summarybox}[n次方根]
		一般地，如果 \(x^n=a\)，那么 \(x\) 叫做 \(a\) 的\textbf{\(n\) 次方根}，其中 \(n>1\)，且 \(n\in \mathbf{N}^{*}\).
	\end{summarybox}
	\vspace{1em}
	\begin{columns}
		\column{0.48\textwidth}
		\begin{summarybox}[当根指数 n 是奇数时]
			正数的 \(n\) 次方根是一个正数，负数的 \(n\) 次方根是一个负数. 这时，\(a\) 的 \(n\) 次方根用符号 \(\sqrt[n]{a}\) 表示.
		\end{summarybox}
		\column{0.48\textwidth}
		\begin{summarybox}[当根指数 n 是偶数时]
			正数的 \(n\) 次方根有两个，这两个数互为相反数. 正的 \(n\) 次方根用 \(\sqrt[n]{a}\) 表示，负的用 \(-\sqrt[n]{a}\) 表示，合并写成 \(\pm\sqrt[n]{a}\;(a>0)\).
		\end{summarybox}
	\end{columns}
\end{frame}

\begin{frame}{几个约定与根式}
	\begin{summarybox}[注意]
		\begin{itemize}
			\item 负数没有偶次方根.
			\item \(0\) 的任何次方根都是 \(0\)，记作 \(\sqrt[n]{0}=0\).
		\end{itemize}
	\end{summarybox}
	\vspace{1em}
	\pause
	例子：
	\[
		\sqrt[5]{32}=2,\quad \sqrt[5]{-32}=-2,\quad \sqrt[3]{a^6}=a^2.
	\]
	\[
		\sqrt[4]{16}=2,\quad -\sqrt[4]{16}=-2,\quad \pm\sqrt[4]{16}=\pm 2.
	\]
	\vspace{-1em}
	\pause
	\begin{summarybox-answer}
		式子 \(\sqrt[n]{a}\) 叫做\textbf{根式（radical）}，这里 \(n\) 叫做\textbf{根指数}，\(a\) 叫做\textbf{被开方数}.
	\end{summarybox-answer}
\end{frame}

\begin{frame}{根式的两个重要性质}
	根据 \(n\) 次方根的意义，可得
	\[
		\left(\sqrt[n]{a}\right)^n=a.
	\]
	例如，\((\sqrt{5})^2=5\)，\((\sqrt[5]{-3})^5=-3\).
	\vspace{1em}
	\pause
	\begin{summarybox}[探究：根式开方规律]
	\hspace{2em}\(\sqrt[n]{a^n}\) 表示$a^n$的$n$次方根，$\sqrt[n]{a^n} = a$一定成立吗？如果不一定成立，那么$\sqrt[n]{a^n}$等于什么？	
	\end{summarybox}
\end{frame}


\begin{frame}{根式的两个重要性质}
	
	\begin{summarybox-answer}
		\hspace{2em}当 \(n\) 为奇数时，\(\sqrt[n]{a^n}=a\)；
		\[
		\text{当 } n \text{ 为偶数时，}\quad \sqrt[n]{a^n}=|a|=\begin{cases} a, & a\geqslant 0,\\ -a, & a<0.\end{cases}
		\]
	\end{summarybox-answer}
\end{frame}

\begin{frame}{例1：求下列各式的值}
	\begin{exam}
		求下列各式的值：
		\begin{enumerate}[(1)]
			\item \(\sqrt[3]{(-8)^3}\)；
			\item \(\sqrt{(-10)^2}\)；
			\item \(\sqrt[4]{(3-\pi)^4}\)；
			\item \(\sqrt{(a-b)^2}\).
		\end{enumerate}
	\end{exam}
\end{frame}

\begin{frame}{例1 解答}
	\begin{solve}
		\begin{enumerate}[(1)]
			\item \(\sqrt[3]{(-8)^3}=-8\)；
			\item \(\sqrt{(-10)^2}=|-10|=10\)；
			\item \(\sqrt[4]{(3-\pi)^4}=|3-\pi|=\pi-3\)；
			\item \(\sqrt{(a-b)^2}=|a-b|\)，当 \(a\geqslant b\) 时为 \(a-b\)，当 \(a<b\) 时为 \(b-a\).
		\end{enumerate}
	\end{solve}
	\vspace{1em}
	\begin{summarybox-answer}
		\textbf{关键：}偶次根式开方后要加绝对值，再根据符号去掉绝对值.
	\end{summarybox-answer}
\end{frame}

\section{分数指数幂}

\begin{frame}{根式化为分数指数幂}
	根据 \(n\) 次方根的定义和数的运算：
	\[
		\sqrt[5]{a^{10}}=\sqrt[5]{(a^2)^5}=a^2=a^{\frac{10}{5}}\quad (a>0),
	\]
	\[
		\sqrt[4]{a^{12}}=\sqrt[4]{(a^3)^4}=a^3=a^{\frac{12}{4}}\quad (a>0).
	\]
	\vspace{0em}
	\pause
	\begin{summarybox-answer}
		当根式的被开方数（看成幂的形式）的指数能被根指数整除时，根式可以表示为分数指数幂的形式.
	\end{summarybox-answer}
\end{frame}

\begin{frame}{分数指数幂的规定}
	\begin{summarybox}[正数的正分数指数幂]
		\[
			a^{\frac{m}{n}}=\sqrt[n]{a^m}\quad (a>0,\ m,n\in \mathbf{N}^{*},\ n>1).
		\]
	\end{summarybox}
	
	\begin{summarybox}[正数的负分数指数幂]
		\[
			a^{-\frac{m}{n}}=\frac{1}{a^{\frac{m}{n}}}=\frac{1}{\sqrt[n]{a^m}}\quad (a>0,\ m,n\in \mathbf{N}^{*},\ n>1).
		\]
	\end{summarybox}
	

	\begin{summarybox}[0的分数指数幂]
		\(0\) 的正分数指数幂等于 \(0\)，\(0\) 的负分数指数幂\textbf{没有意义}.
	\end{summarybox}
\end{frame}

\begin{frame}{有理数指数幂的运算性质}
	规定了分数指数幂以后，幂 \(a^x\) 中指数 \(x\) 的取值范围就从整数拓展到了有理数. 整数指数幂的运算性质对有理数指数幂同样适用：
	\vspace{1em}
	\begin{summarybox}[运算性质（底数大于0）]
		\begin{enumerate}[(1)]
			\item \(a^r\cdot a^s=a^{r+s}\)；
			\item \((a^r)^s=a^{rs}\)；
			\item \((ab)^r=a^r b^r\quad (a>0,\ b>0,\ r\in \mathbf{Q})\).
		\end{enumerate}
	\end{summarybox}
\end{frame}

\begin{frame}{例2：求值}
	\begin{exam}
		求值：
		\begin{enumerate}[(1)]
			\item \(8^{\frac{2}{3}}\)；
			\item \(\left(\dfrac{16}{81}\right)^{-\frac{3}{4}}\).
		\end{enumerate}
	\end{exam}
	\pause
	\begin{solve}
		\begin{enumerate}[(1)]
			\item \(8^{\frac{2}{3}}=(2^3)^{\frac{2}{3}}=2^{3\times \frac{2}{3}}=2^2=4\)；
			\item \(\left(\dfrac{16}{81}\right)^{-\frac{3}{4}}=\left(\dfrac{81}{16}\right)^{\frac{3}{4}}=\left[\left(\dfrac{3}{2}\right)^4\right]^{\frac{3}{4}}=\left(\dfrac{3}{2}\right)^3=\dfrac{27}{8}\).
		\end{enumerate}
	\end{solve}
\end{frame}

\begin{frame}{例3：用分数指数幂表示并计算}
	\begin{exam}
		用分数指数幂的形式表示并计算下列各式（其中 \(a>0\)）：
		\begin{enumerate}[(1)]
			\item \(a^2\cdot \sqrt[3]{a^2}\)；
			\item \(\sqrt{a\sqrt[3]{a}}\).
		\end{enumerate}
	\end{exam}
	\pause
	\begin{solve}
		\begin{enumerate}[(1)]
			\item \(a^2\cdot \sqrt[3]{a^2}=a^2 a^{\frac{2}{3}}=a^{2+\frac{2}{3}}=a^{\frac{8}{3}}\)；
			\item \(\sqrt{a\sqrt[3]{a}}=(a\cdot a^{\frac{1}{3}})^{\frac{1}{2}}=(a^{\frac{4}{3}})^{\frac{1}{2}}=a^{\frac{2}{3}}\).
		\end{enumerate}
	\end{solve}
\end{frame}

\begin{frame}{}
	\begin{exam}
		计算下列各式（式中字母均是正数）：
		\begin{enumerate}[(1)]
			\item \(\left(2a^{\frac{2}{3}}b^{\frac{1}{2}}\right)\left(-6a^{\frac{1}{2}}b^{\frac{1}{3}}\right)\div\left(-3a^{\frac{5}{6}}b^{\frac{1}{6}}\right)\)；
			\item \(\left(m^{\frac{1}{4}}n^{-\frac{3}{8}}\right)^8\);
			\item $(\sqrt[3]{a^2} - \sqrt{a^3}) \div \sqrt[4]{a^2}$.
		\end{enumerate}
	\end{exam}
	\pause
	\begin{solve}
		\begin{enumerate}[(1)]
			\item 原式 \(=[2\times(-6)\div(-3)]\,a^{\frac{2}{3}+\frac{1}{2}-\frac{5}{6}}b^{\frac{1}{2}+\frac{1}{3}-\frac{1}{6}}=4ab^{0}=4a\)；
			\item 原式 \(=(m^{\frac{1}{4}})^8(n^{-\frac{3}{8}})^8=m^2 n^{-3}=\dfrac{m^2}{n^3}\).
			\item $(\sqrt[3]{a^2} - \sqrt{a^3}) \div \sqrt[4]{a^2}=(a^{\frac{2}{3}} - a^{\frac{3}{2}}) \div a^{\frac{1}{2}} = a^{\frac{2}{3} \div a^{\frac{1}{2}} - a^{\frac{3}{2}}} \div a^{\frac{1}{2}} = a^{\frac{1}{6}} -a = \sqrt[6]{a} -a.$
		\end{enumerate}
	\end{solve}
\end{frame}

\begin{frame}{布置作业}
	\begin{summarybox}[课后作业]
		\begin{enumerate}
			\item 教材 P107 练习 1、2、3
			\item 教材 P109 习题 4.1：1---5
			
		\end{enumerate}
	\end{summarybox}
\end{frame}

%\section{无理数指数幂}
%
%\begin{frame}{无理数指数幂}
%	\begin{summarybox}[探究]
%		根据 \(\sqrt{2}\) 的不足近似值 \(x\) 和过剩近似值 \(y\)，利用计算工具计算相应的 \(5^x\)，\(5^y\) 的近似值并填入表中，观察它们的变化趋势.
%	\end{summarybox}
%	\vspace{0.5em}
%	\begin{center}
%		\small
%		\begin{tabular}{|c|c|c|c|}
%			\hline
%			\(\sqrt{2}\) 的不足近似值 \(x\) & \(5^x\) 的近似值 & \(\sqrt{2}\) 的过剩近似值 \(y\) & \(5^y\) 的近似值 \\ \hline
%			1.4 & 9.518 & 1.5 & 11.118 \\ \hline
%			1.41 & 9.672 & 1.42 & 9.829 \\ \hline
%			1.414 & 9.735 & 1.415 & 9.750 \\ \hline
%			\(\cdots\) & \(\cdots\) & \(\cdots\) & \(\cdots\) \\ \hline
%		\end{tabular}
%	\end{center}
%\end{frame}
%
%\begin{frame}{实数指数幂的运算性质}
%	可以发现，当 \(\sqrt{2}\) 的不足近似值 \(x\) 和过剩近似值 \(y\) 逐渐逼近 \(\sqrt{2}\) 时，\(5^x\) 和 \(5^y\) 都趋向于同一个数，这个数就是 \(5^{\sqrt{2}}\)，它是一个\textbf{确定的实数}.
%	\vspace{1em}
%	\pause
%	\begin{summarybox}[无理数指数幂]
%		一般地，无理数指数幂 \(a^{\alpha}\)（\(a>0\)，\(\alpha\) 为无理数）是一个\textbf{确定的实数}. 这样，我们就将指数幂 \(a^x\)（\(a>0\)）中指数 \(x\) 的取值范围从整数逐步拓展到了\textbf{实数}.
%	\end{summarybox}
%	\vspace{1em}
%	\pause
%	\begin{summarybox-answer}
%		整数指数幂的运算性质也适用于实数指数幂，即对任意实数 \(r,s\)，均有
%		\[
%			a^r a^s=a^{r+s},\quad (a^r)^s=a^{rs},\quad (ab)^r=a^r b^r\quad (a>0,b>0,r,s\in \mathbf{R}).
%		\]
%	\end{summarybox-answer}
%\end{frame}
%
%\section{课堂练习}
%
%\begin{frame}{练习 1：根式与分数指数幂互化}
%	\begin{question}
%		用根式的形式表示下列各式（\(a>0\)）：
%		\[
%			a^{\frac{1}{5}},\qquad a^{\frac{3}{4}},\qquad a^{-\frac{3}{5}},\qquad a^{-\frac{2}{3}}.
%		\]
%	\end{question}
%	\pause
%	\begin{solve}
%		\(a^{\frac{1}{5}}=\sqrt[5]{a}\)；\quad \(a^{\frac{3}{4}}=\sqrt[4]{a^3}\)；\\[4pt]
%		\(a^{-\frac{3}{5}}=\dfrac{1}{\sqrt[5]{a^3}}\)；\quad \(a^{-\frac{2}{3}}=\dfrac{1}{\sqrt[3]{a^2}}\).
%	\end{solve}
%\end{frame}
%
%\begin{frame}{练习 2：计算}
%	\begin{question}
%		计算下列各式：
%		\begin{enumerate}[(1)]
%			\item \(\left(\dfrac{36}{49}\right)^{\frac{3}{2}}\)；
%			\item \(2\sqrt{3}\times 3\sqrt[3]{1.5}\times \sqrt[6]{12}\).
%		\end{enumerate}
%	\end{question}
%	\pause
%	\begin{solve}
%		\begin{enumerate}[(1)]
%			\item 原式 \(=\left[\left(\dfrac{6}{7}\right)^2\right]^{\frac{3}{2}}=\left(\dfrac{6}{7}\right)^3=\dfrac{216}{343}\)；
%			\item 原式 \(=2\cdot 3^{\frac{1}{2}}\cdot 3\cdot \left(\dfrac{3}{2}\right)^{\frac{1}{3}}\cdot (12)^{\frac{1}{6}}=6\cdot 3^{\frac{1}{2}+\frac{1}{3}+\frac{1}{6}}\cdot 2^{-\frac{1}{3}+\frac{1}{3}}=6\cdot 3=18\).
%		\end{enumerate}
%	\end{solve}
%\end{frame}



\end{document}
